\documentclass[12pt,a4paper]{amsart}

\usepackage{amsmath,amssymb,amsthm}
\usepackage{mathtools}
\usepackage{hyperref}
\usepackage{geometry}
\usepackage[ngerman]{babel}
\geometry{margin=2.5cm}

% -------------------------------------------------------
\newtheorem{theorem}{Satz}[section]
\newtheorem{lemma}[theorem]{Lemma}
\newtheorem{corollary}[theorem]{Korollar}
\newtheorem{proposition}[theorem]{Proposition}
\theoremstyle{definition}
\newtheorem{definition}[theorem]{Definition}
\newtheorem{remark}[theorem]{Bemerkung}

\newcommand{\Z}{\mathbb{Z}}
\newcommand{\euler}{\varphi}

% -------------------------------------------------------
\begin{document}
% -------------------------------------------------------

\title{Jede Lösung von Lehmers Totient-Problem ist quadratfrei:\\ Ein kurzer Beweis}

\author{Kurt Ingwer}
\address{Specialist Mathematics Project, Version Build 44}
\email{specialist-maths@localhost}
\date{\today}

\begin{abstract}
Lehmers Totient-Problem (1932) fragt, ob die Bedingung $\euler(n) \mid (n-1)$
die Primalität von~$n$ erzwingt. Eine zusammengesetzte Lösung würde als
\emph{Lehmer-Zahl} bezeichnet. Wir geben einen kurzen, elementaren Beweis, dass
jede Lehmer-Zahl quadratfrei sein muss. Das Argument besteht aus zwei Schritten:
Falls $p^2 \mid n$, dann $p \mid \euler(n)$; zusammen mit $\euler(n) \mid (n-1)$
folgt $p \mid (n-1)$; aber auch $p \mid n$, also $p \mid 1$, was unmöglich ist.
Dieses Ergebnis ist grundlegend für das Studium von Lehmers Problem und impliziert,
dass jedes potenzielle Gegenbeispiel ein Produkt verschiedener Primzahlen ist.
\end{abstract}

\maketitle

% -------------------------------------------------------
\section{Einleitung}
% -------------------------------------------------------

Im Jahre 1932 fragte Lehmer \cite{Lehmer1932}, ob die Kongruenz
$\euler(n) \mid (n-1)$ (wobei $\euler$ die Eulersche Totient-Funktion ist)
für eine zusammengesetzte ganze Zahl~$n$ gelten kann. Primzahlen erfüllen stets
$\euler(p) = p-1 \mid p-1$; das Problem ist, ob dies die Primzahlen charakterisiert.

\begin{definition}
Eine zusammengesetzte ganze Zahl~$n$ mit $\euler(n) \mid (n-1)$ heißt
\emph{Lehmer-Zahl} oder \emph{Lehmer-Pseudoprim}.
\end{definition}

Keine Lehmer-Zahl wurde jemals gefunden. Cohen und Hagis \cite{Cohen1980} bewiesen,
dass eine solche Zahl mindestens~$13$ Primfaktoren haben muss; Banks und Luca
\cite{Banks2009} verbesserten dies auf~$15$.

Die vorliegende Notiz beweist die grundlegendste Struktureigenschaft: die Quadratfreiheit.

% -------------------------------------------------------
\section{Hauptresultat}
% -------------------------------------------------------

\begin{theorem}
\label{thm:qfrei}
Jede Lehmer-Zahl ist quadratfrei.
\end{theorem}

\begin{proof}
Angenommen, $n$ ist eine Lehmer-Zahl und $p^2 \mid n$ für eine Primzahl~$p$.

Da $p^2 \mid n$, gilt $p \mid \euler(n)$.
(Zur Erinnerung: Falls $p^a \| n$ mit $a \ge 2$, dann teilt $p^{a-1}(p-1)$ die Zahl $\euler(n)$,
insbesondere gilt also $p \mid \euler(n)$; schon für $a = 2$ liefert dies $p(p-1) \mid \euler(n)$.)

Aus $\euler(n) \mid (n-1)$ und $p \mid \euler(n)$ folgt:
\[
  p \mid \euler(n) \mid (n-1), \quad \text{also} \quad p \mid (n-1).
\]
Aber auch $p \mid n$ (da $p^2 \mid n$), daher
\[
  p \mid n - (n-1) = 1.
\]
Dies ist unmöglich für $p \ge 2$ — Widerspruch.
Folglich erfüllt keine Primzahl~$p$ die Bedingung $p^2 \mid n$, d.\,h.\ $n$ ist quadratfrei.
\end{proof}

\begin{remark}
Dasselbe Argument gilt für die allgemeinere Bedingung $M \cdot \euler(n) = n - 1$
für jedes feste positive ganze Zahl~$M$ (eine Verallgemeinerung, die von Hagis \cite{Hagis1988}
untersucht wurde). Quadratfreiheit folgt identisch.
\end{remark}

% -------------------------------------------------------
\section{Folgerungen}
% -------------------------------------------------------

\begin{corollary}
Jede Lehmer-Zahl hat die Form $n = p_1 p_2 \cdots p_k$ mit paarweise verschiedenen
Primzahlen $p_1 < p_2 < \cdots < p_k$ und $k \ge 2$.
\end{corollary}

\begin{corollary}
Für eine Lehmer-Zahl $n = p_1 \cdots p_k$ ist der Totient
$\euler(n) = (p_1-1)(p_2-1)\cdots(p_k-1)$, und die Lehmer-Bedingung lautet
\[
  (p_1-1)(p_2-1)\cdots(p_k-1) \mid p_1 p_2 \cdots p_k - 1.
\]
\end{corollary}

Diese explizite Form der Bedingung ist die Grundlage für alle weitere Analyse von
Lehmer-Zahlen, einschließlich der Nachweise, dass kein Semiprim \cite{Ingwer2026b}
und keine Dreiprimzahl \cite{Ingwer2026lehmer3} eine Lehmer-Zahl sein kann.

% -------------------------------------------------------
\section{Vergleich mit dem Giuga-Fall}
% -------------------------------------------------------

Die parallele Struktur mit Giuga-Pseudoprimen ist auffällig. Für beide Probleme:
\begin{itemize}
  \item Quadratfreiheit wird bewiesen, indem gezeigt wird, dass $p^2 \mid n$ zu $p \mid 1$ führt.
  \item Für Giuga: $p^2 \mid n \Rightarrow p \mid \tfrac{n}{p} \Rightarrow \tfrac{n}{p}-1 \equiv -1 \pmod{p} \Rightarrow p \mid -1$.
  \item Für Lehmer: $p^2 \mid n \Rightarrow p \mid \euler(n) \Rightarrow p \mid (n-1) \Rightarrow p \mid 1$.
\end{itemize}
Beide nutzen die fundamentale Unverträglichkeit zwischen $p \mid n$ und $p \mid (n-1)$.

% -------------------------------------------------------
\begin{thebibliography}{10}

\bibitem{Lehmer1932}
D.~H.~Lehmer,
\emph{On Euler's totient function},
Bull.\ Amer.\ Math.\ Soc.\ \textbf{38} (1932), 745--751.

\bibitem{Cohen1980}
G.~L.~Cohen und P.~Hagis,
\emph{On the number of prime factors of $n$ if $\phi(n) \mid (n-1)$},
Nieuw Arch.\ Wisk.\ (3) \textbf{28} (1980), 177--185.

\bibitem{Hagis1988}
P.~Hagis,
\emph{On the equation $M \cdot \phi(n) = n - 1$},
Nieuw Arch.\ Wisk.\ (4) \textbf{6} (1988), 255--261.

\bibitem{Banks2009}
W.~D.~Banks und F.~Luca,
\emph{Composite integers $n$ for which $\phi(n) \mid n - 1$},
Acta Math.\ Sin.\ (Engl.\ Ser.) \textbf{25} (2009), 1703--1710.

\bibitem{Ingwer2026b}
K.~Ingwer,
\emph{Kein Semiprim erfüllt Lehmers Totient-Bedingung},
Specialist Mathematics Project (2026).

\bibitem{Ingwer2026lehmer3}
K.~Ingwer,
\emph{Lehmers Totient-Problem für Produkte dreier Primzahlen},
Specialist Mathematics Project (2026).

\end{thebibliography}

\end{document}
